% ID: SEL-D0082-OB-E04
% Tanda S082.
\textbf{a)} Contrastamos $H_0:p\le0{,}05$ frente a $H_1:p>0{,}05$. La alternativa requiere una región crítica en la cola derecha. En la frontera $p_0=0{,}05$, usamos la aproximación normal del estadístico
\[
 Z=\frac{\hat p-0{,}05}{\sqrt{0{,}05\cdot0{,}95/300}}.
\]
Las frecuencias esperadas bajo esa frontera son 15 defectuosos y 285 no defectuosos. Para $\alpha=0{,}05$, el valor crítico es $z_{0{,}95}\approx1{,}64485$, por lo que la región crítica es
\[
 \boxed{Z>1{,}64485},
 \qquad\text{equivalentemente}\qquad\hat p>0{,}070697.
\]
Para proporciones inferiores a $0{,}05$, el rechazo por proporciones observadas altas es menos probable. Se utiliza la frontera para fijar el contraste.

\textbf{b)} La proporción observada es $\hat p=35/300\approx0{,}116667$. Por tanto,
\[
 Z_{\mathrm{obs}}\approx5{,}2981>1{,}64485.
\]
Rechazamos $H_0$ al $5\%$. Los datos aportan evidencia de que la proporción de artículos defectuosos ha aumentado por encima del $5\%$. El valor $p$ de la aproximación normal es aproximadamente $5{,}85\cdot10^{-8}$.
